\textbf{Phần 1. Trắc nghiệm lựa chọn}
\setcounterpageref{ex}{0}

\begin{ex} %[5H9N55]
Hình lập phương có cạnh $a$ thì diện tích xung quanh gấp mấy lần diện tích một mặt?
\choice
{\True $4$ lần}
{$2$ lần}
{$6$ lần}
{$8$ lần}
\loigiai{
$S_{xq} = a \times a \times 4$ nên gấp $4$ lần diện tích một mặt.\\
\textbf{Đáp án của câu hỏi là A}
}
\end{ex}

\begin{ex} %[5H9N55]
Một hình hộp chữ nhật có diện tích xung quanh là $120\text{ cm}^2$ và diện tích hai đáy là $80\text{ cm}^2$. Diện tích toàn phần của hình hộp là
\choice
{$160\text{ cm}^2$}
{$240\text{ cm}^2$}
{\True $200\text{ cm}^2$}
{$280\text{ cm}^2$}
\loigiai{
$S_{tp} = S_{xq} + S_{2đáy} = 120 + 80 = 200\text{ (cm}^2\text{)}$.\\
\textbf{Đáp án của câu hỏi là C}
}
\end{ex}

\begin{ex} %[5H9H55]
Một hình lập phương có diện tích toàn phần là $150\text{ dm}^2$. Thể tích của nó là
\choice
{$100\text{ dm}^3$}
{$150\text{ dm}^3$}
{$250\text{ dm}^3$}
{\True $125\text{ dm}^3$}
\loigiai{
Diện tích 1 mặt: $150 : 6 = 25\text{ (dm}^2\text{)}$. Cạnh là $5\text{ dm}$.\\
$V = 5 \times 5 \times 5 = 125\text{ (dm}^3\text{)}$.\\
\textbf{Đáp án của câu hỏi là D}
}
\end{ex}

\begin{ex} %[5H9H55]
Một bể nước hình hộp chữ nhật có chiều dài $2\text{ m}$, chiều rộng $1{,}5\text{ m}$, chiều cao $1\text{ m}$. Bể chứa được tối đa số lít nước là
\choice
{$300\text{ l}$}
{\True $3\,000\text{ l}$}
{$30\,000\text{ l}$}
{$30\text{ l}$}
\loigiai{
$V = 2 \times 1{,}5 \times 1 = 3\text{ m}^3 = 3\,000\text{ l}$.\\
\textbf{Đáp án của câu hỏi là B}
}
\end{ex}

\begin{ex} %[5H9H55]
Diện tích xung quanh của một hình lập phương là $64\text{ cm}^2$. Diện tích toàn phần của nó là
\choice
{$80\text{ cm}^2$}
{$120\text{ cm}^2$}
{$144\text{ cm}^2$}
{\True $96\text{ cm}^2$}
\loigiai{
Diện tích 1 mặt: $64 : 4 = 16\text{ (cm}^2\text{)}$.\\
$S_{tp} = 16 \times 6 = 96\text{ (cm}^2\text{)}$.\\
\textbf{Đáp án của câu hỏi là D}
}
\end{ex}

\begin{ex} %[5H9H55]
Một khối gỗ hình hộp chữ nhật có thể tích $0{,}72\text{ m}^3$, chiều dài $1{,}2\text{ m}$, chiều rộng $0{,}8\text{ m}$. Chiều cao của khối gỗ là
\choice
{$0{,}6\text{ m}$}
{\True $0{,}75\text{ m}$}
{$0{,}8\text{ m}$}
{$0{,}9\text{ m}$}
\loigiai{
Chiều cao: $0{,}72 : (1{,}2 \times 0{,}8) = 0{,}72 : 0{,}96 = 0{,}75\text{ (m)}$.\\
\textbf{Đáp án của câu hỏi là B}
}
\end{ex}

\begin{ex} %[5H9V55]
Người ta quét sơn toàn bộ mặt ngoài của một chiếc thùng gỗ có nắp dạng hình hộp chữ nhật với chiều dài $1{,}6\text{ m}$, chiều rộng $1{,}2\text{ m}$, chiều cao $0{,}8\text{ m}$. Diện tích cần quét sơn là
\choice
{$4{,}48\text{ m}^2$}
{$6{,}4\text{ m}^2$}
{\True $8{,}32\text{ m}^2$}
{$7{,}28\text{ m}^2$}
\loigiai{
$S_{xq} = (1{,}6 + 1{,}2) \times 2 \times 0{,}8 = 4{,}48\text{ (m}^2\text{)}$.\\
$S_{2đáy} = 2 \times 1{,}6 \times 1{,}2 = 3{,}84\text{ (m}^2\text{)}$.\\
$S_{tp} = 4{,}48 + 3{,}84 = 8{,}32\text{ (m}^2\text{)}$.\\
\textbf{Đáp án của câu hỏi là C}
}
\end{ex}

\begin{ex} %[5H9V55]
Một hình hộp chữ nhật và một hình lập phương có cùng thể tích. Biết hình hộp chữ nhật có kích thước dài $8\text{ cm}$, rộng $4\text{ cm}$, cao $2\text{ cm}$. Cạnh của hình lập phương là
\choice
{$2\text{ cm}$}
{\True $4\text{ cm}$}
{$6\text{ cm}$}
{$8\text{ cm}$}
\loigiai{
$V = 8 \times 4 \times 2 = 64\text{ (cm}^3\text{)}$.\\
Vì $4 \times 4 \times 4 = 64$ nên cạnh hình lập phương là $4\text{ cm}$.\\
\textbf{Đáp án của câu hỏi là B}
}
\end{ex}

\begin{ex} %[5H9V55]
Một bể nước hình hộp chữ nhật có diện tích đáy $4\text{ m}^2$, chiều cao $1{,}5\text{ m}$. Hiện tại bể đang có $4\,200\text{ l}$ nước. Lượng nước hiện có chiếm bao nhiêu phần trăm thể tích bể?
\choice
{\True $70\%$}
{$60\%$}
{$80\%$}
{$75\%$}
\loigiai{
Thể tích bể: $4 \times 1{,}5 = 6\text{ m}^3 = 6\,000\text{ l}$.\\
Tỉ số phần trăm: $4\,200 : 6\,000 = 0{,}7 = 70\%$.\\
\textbf{Đáp án của câu hỏi là A}
}
\end{ex}

\begin{ex} %[5H9V55]
Một hình lập phương có cạnh $a$. Nếu tăng cạnh $a$ lên $10\%$ thì diện tích toàn phần của nó tăng thêm bao nhiêu phần trăm?
\choice
{$10\%$}
{$20\%$}
{\True $21\%$}
{$33{,}1\%$}
\loigiai{
Cạnh mới: $1{,}1a$.\\
Diện tích mới: $(1{,}1a)^2 \times 6 = 1{,}21 \times (a^2 \times 6) = 1{,}21 S_{tp} = 121\% S_{tp}$.\\
Tăng thêm: $121\% - 100\% = 21\%$.\\
\textbf{Đáp án của câu hỏi là C}
}
\end{ex}

\begin{ex} %[5H9V55]
Một cái thùng tôn không nắp dạng hình lập phương có diện tích tôn là $125\text{ dm}^2$. Thể tích của cái thùng đó là
\choice
{$100\text{ dm}^3$}
{$150\text{ dm}^3$}
{$625\text{ dm}^3$}
{\True $125\text{ dm}^3$}
\loigiai{
Thùng không nắp có 5 mặt $\Rightarrow$ Diện tích 1 mặt: $125 : 5 = 25\text{ (dm}^2\text{)}$.\\
Cạnh là $5\text{ dm} \Rightarrow V = 5 \times 5 \times 5 = 125\text{ (dm}^3\text{)}$.\\
\textbf{Đáp án của câu hỏi là D}
}
\end{ex}

\begin{ex} %[5H9C55]
Một bể cá hình hộp chữ nhật dài $1\text{ m}$, rộng $0{,}6\text{ m}$, cao $0{,}5\text{ m}$. Bể đang chứa nước cao $0{,}3\text{ m}$. Người ta thả vào bể một khối kim loại hình lập phương cạnh $20\text{ cm}$ thì mực nước trong bể dâng lên cao bao nhiêu mét (khối kim loại chìm hoàn toàn)?
\choice
{\True $0{,}313\text{ m}$}
{$0{,}32\text{ m}$}
{$0{,}35\text{ m}$}
{$0{,}38\text{ m}$}
\loigiai{
Thể tích khối lập phương: $0{,}2 \times 0{,}2 \times 0{,}2 = 0{,}008\text{ (m}^3\text{)}$.\\
Diện tích đáy bể: $1 \times 0{,}6 = 0{,}6\text{ (m}^2\text{)}$.\\
Chiều cao nước dâng thêm: $0{,}008 : 0{,}6 = \dfrac{0{,}008}{0{,}6} = \dfrac{8}{600} = \dfrac{1}{75}\text{ m} \approx 0{,}0133\text{ m}$.\\
Mực nước mới: $0{,}3 + 0{,}0133 = 0{,}3133\text{ m} \approx 0{,}313\text{ m}$.\\
\textbf{Đáp án của câu hỏi là A}
}
\end{ex}

\textbf{Phần 2. Trắc nghiệm đúng sai}
\setcounterpageref{ex}{0}

\begin{ex} %[5H9H55]
So sánh hình hộp chữ nhật có kích thước $6\text{ cm} \times 4\text{ cm} \times 3\text{ cm}$ và hình lập phương có cạnh $4\text{ cm}$. Xét tính đúng sai:
\choiceTF
{\True Diện tích toàn phần của hình hộp chữ nhật là $108\text{ cm}^2$.}
{\True Diện tích toàn phần của hình lập phương là $96\text{ cm}^2$.}
{Thể tích hình hộp chữ nhật lớn hơn thể tích hình lập phương.}
{\True Thể tích hình lập phương bé hơn thể tích hình hộp chữ nhật ($64\text{ cm}^3 < 72\text{ cm}^3$).}
\loigiai{
\begin{itemize}
    \item $S_{tp(HHCN)} = 2 \times (24 + 18 + 12) = 108\text{ cm}^2$. Đúng.
    \item $S_{tp(HLP)} = 4 \times 4 \times 6 = 96\text{ cm}^2$. Đúng.
    \item $V_{HHCN} = 6 \times 4 \times 3 = 72\text{ cm}^3$; $V_{HLP} = 64\text{ cm}^3$. HHCN lớn hơn HLP: Đúng.
    \item Thể tích HLP bé hơn HHCN: Đúng.
\end{itemize}
\begin{itemize}
    \item a) \textbf{Đúng}.
    \item b) \textbf{Đúng}.
    \item c) \textbf{Đúng}.
    \item d) \textbf{Đúng}.
\end{itemize}
}
\end{ex}

\begin{ex} %[5H9V55]
Một cái bể nước hình hộp chữ nhật dài $2{,}5\text{ m}$, rộng $1{,}6\text{ m}$, cao $1{,}2\text{ m}$. Xét tính đúng sai:
\choiceTF
{\True Diện tích xung quanh lòng bể là $9{,}84\text{ m}^2$.}
{\True Diện tích đáy bể là $4\text{ m}^2$.}
{Dung tích toàn bộ bể là $4\,800\text{ l}$.}
{\True Khi bể chứa $3\,600\text{ l}$ nước thì mực nước cao $0{,}9\text{ m}$.}
\loigiai{
\begin{itemize}
    \item $S_{xq} = (2{,}5 + 1{,}6) \times 2 \times 1{,}2 = 9{,}84\text{ m}^2$. Đúng.
    \item $S_{đáy} = 2{,}5 \times 1{,}6 = 4\text{ m}^2$. Đúng.
    \item Dung tích: $4 \times 1{,}2 = 4{,}8\text{ m}^3 = 4800\text{ l}$. Đúng.
    \item Chiều cao khi chứa 3600 l: $3{,}6 : 4 = 0{,}9\text{ m}$. Đúng.
\end{itemize}
\begin{itemize}
    \item a) \textbf{Đúng}.
    \item b) \textbf{Đúng}.
    \item c) \textbf{Đúng}.
    \item d) \textbf{Đúng}.
\end{itemize}
}
\end{ex}

\textbf{Phần 3. Trả lời ngắn}
\setcounterpageref{ex}{0}

\begin{ex} %[5H9H55]
Một hình lập phương có cạnh $11\text{ cm}$. Tính thể tích theo xăng-ti-mét khối.
\shortans{1331}
\loigiai{
\textbf{Đáp số:} $1331$\\
$11 \times 11 \times 11 = 1\,331\text{ cm}^3$.
}
\end{ex}

\begin{ex} %[5H9H55]
Một hình hộp chữ nhật có diện tích xung quanh $180\text{ cm}^2$, chiều cao $6\text{ cm}$. Tính chu vi đáy theo xăng-ti-mét.
\shortans{30}
\loigiai{
\textbf{Đáp số:} $30$\\
$180 : 6 = 30\text{ cm}$.
}
\end{ex}

\begin{ex} %[5H9V55]
Tính diện tích toàn phần của hình hộp chữ nhật có kích thước dài $15\text{ dm}$, rộng $8\text{ dm}$, cao $10\text{ dm}$ theo đề-xi-mét vuông.
\shortans{700}
\loigiai{
\textbf{Đáp số:} $700$\\
$S_{xq} = (15 + 8) \times 2 \times 10 = 460\text{ dm}^2$.\\
$S_{2đáy} = 2 \times 15 \times 8 = 240\text{ dm}^2$.\\
$S_{tp} = 460 + 240 = 700\text{ dm}^2$.
}
\end{ex}

\begin{ex} %[5H9V55]
Một bể bơi hình hộp chữ nhật dài $18\text{ m}$, rộng $8\text{ m}$, sâu $1{,}5\text{ m}$. Bể đang chứa $144\text{ m}^3$ nước. Hỏi cần bơm thêm bao nhiêu mét khối nước nữa để đầy bể?
\shortans{72}
\loigiai{
\textbf{Đáp số:} $72$\\
Thể tích bể: $18 \times 8 \times 1{,}5 = 216\text{ m}^3$.\\
Cần thêm: $216 - 144 = 72\text{ m}^3$.
}
\end{ex}

\begin{ex} %[5H9V55]
Một hình lập phương có diện tích toàn phần là $216\text{ cm}^2$. Tính thể tích của hình lập phương đó theo xăng-ti-mét khối.
\shortans{216}
\loigiai{
\textbf{Đáp số:} $216$\\
Diện tích 1 mặt: $216 : 6 = 36\text{ cm}^2 \Rightarrow$ Cạnh $6\text{ cm}$.\\
$V = 6 \times 6 \times 6 = 216\text{ cm}^3$.
}
\end{ex}

\begin{ex} %[5H9C55]
Một bể nước hình hộp chữ nhật dài $2\text{ m}$, rộng $1{,}5\text{ m}$, cao $1{,}2\text{ m}$. Người ta dùng máy bơm công suất $25\text{ l/phút}$ để bơm nước vào bể. Hỏi sau bao nhiêu giờ thì bơm đầy bể?
\shortans{2,4}
\loigiai{
\textbf{Đáp số:} $2{,}4$\\
Thể tích: $2 \times 1{,}5 \times 1{,}2 = 3{,}6\text{ m}^3 = 3\,600\text{ l}$.\\
Thời gian: $3\,600 : 25 = 144\text{ phút} = 2{,}4\text{ giờ}$.
}
\end{ex}
